% assets/w06-berth-flow.svg — 접안은 한 번에 붙이지 않는다. 단계가 셋이다
%
%   bash _tools/tikz2svg.sh assets/src/w06-berth-flow.tex assets/w06-berth-flow.svg
%
% 이 그림이 말하는 것
%   접안 제어기는 "선석으로 가라" 한 줄이 아니다. 목표가 단계마다 **바뀐다**.
%   무엇이 바뀌고, 무엇을 보고 다음 단계로 넘어가는지를 한 장에 적은 것이다.
%
%   단계는 한 번 올라가면 내려오지 않는다 (래치). 파도에 밀려 판정 거리를
%   잠깐 넘어서도 앞 단계로 돌아가지 않는다 — 돌아가면 두 단계 사이를 오가며
%   영원히 붙지 못한다.
%
% 판정값은 모델이 쓰는 값 그대로다
%   R_dock = 4 m · 선수각 문턱 5 deg · R_stop = 1.5 m · L_app = 15 m
%   7단(전진)만 둘을 더 쓴다 — R_brake = 8 m · U_berth = 0.2 m/s
%
% 좌표 메모 (mm)
%   가운데 줄 x = 0 · 되돌아가는 선 x = 56 · 왼쪽 설명 x = -118
%   블록 y = -30 / -84 / -138 · 판정 y = -56 / -110 / -164 · 끝 y = -192

\documentclass[border=6pt]{standalone}
% capstone-style.tex — 이 볼트의 TikZ 그림이 공유하는 스타일
%
% 쓰는 법:  % capstone-style.tex — 이 볼트의 TikZ 그림이 공유하는 스타일
%
% 쓰는 법:  % capstone-style.tex — 이 볼트의 TikZ 그림이 공유하는 스타일
%
% 쓰는 법:  \input{capstone-style}   (assets/src/ 안에서)
% 빌드:     bash _tools/tikz2svg.sh assets/src/w02-node-topic.tex assets/w02-node-topic.svg
%
% 왜 TikZ 인가
%   손으로 SVG 좌표를 쓰면 화살촉·선 굵기·글자 크기가 그림마다 미묘하게 달라진다.
%   그것이 "자료가 어설퍼 보인다"의 정체다. 스타일을 한 번만 정하고 상대좌표로
%   배치하면 좌표를 고쳐도 그림이 무너지지 않는다.
%   수식은 본문과 같은 TeX 글꼴로 나온다.
%
% 한글
%   ko.TeX(DVI 경로)를 쓴다. XeLaTeX 은 TikZ 그래픽을 PDF special 로 내보내는데
%   dvisvgm 이 그것을 받으려면 Ghostscript < 10.01 이 필요하다. 이 PC 에는
%   10.07 이 깔려 있어 그 경로가 막힌다. latex -> DVI -> dvisvgm 이 유일하게
%   동작하는 조합이다.

\usepackage[hangul]{kotex}
\usepackage{tikz}
\usepackage{amsmath}
\usepackage{xcolor}
\usetikzlibrary{arrows.meta, positioning, calc, fit, backgrounds, decorations.pathmorphing}

% ---- 색 --------------------------------------------------------------------
% 강조는 하나만. 나머지는 검정. 흑백 인쇄에서도 읽혀야 한다.
\definecolor{ink}{HTML}{1A1A1A}
\definecolor{soft}{HTML}{666666}
\definecolor{accent}{HTML}{7C3AED}
\definecolor{warn}{HTML}{B91C1C}

% 계층 색 — Simulink 모델의 블록 색과 같은 규칙을 쓴다
\definecolor{guid}{HTML}{1D4ED8}    % 유도
\definecolor{ctrl}{HTML}{EA580C}    % 제어
\definecolor{thr}{HTML}{D97706}     % 추진기
\definecolor{plant}{HTML}{16A34A}   % 운동모델
\definecolor{nav}{HTML}{7C3AED}     % 항법 · 신호처리
\definecolor{pathc}{HTML}{0D9488}   % 계획 경로 · 항적
\definecolor{errc}{HTML}{C2410C}    % 오차
\definecolor{flow}{HTML}{0369A1}    % 조류 · 바람

% ---- 선과 화살촉 -----------------------------------------------------------
\tikzset{
  sig/.style   = {ink, line width=0.5pt, -{Stealth[length=4pt, width=3pt]}},
  wire/.style  = {ink, line width=0.5pt},
  fb/.style    = {ink, line width=0.5pt, densely dashed,
                  -{Stealth[length=4pt, width=3pt]}},
  asig/.style  = {accent, line width=0.5pt, densely dashed,
                  -{Stealth[length=4pt, width=3pt]}},
  % 블록 — 흰 바탕, 직각, 검은 테두리. 색은 테두리에만 준다
  blk/.style   = {draw=ink, line width=0.55pt, fill=white,
                  minimum width=18mm, minimum height=8.5mm, inner sep=2pt, font=\small},
  gblk/.style  = {blk, draw=guid},
  cblk/.style  = {blk, draw=ctrl},
  tblk/.style  = {blk, draw=thr},
  pblk/.style  = {blk, draw=plant},
  nblk/.style  = {blk, draw=nav},
  % 합산점 · 분기점
  sum/.style   = {draw=ink, line width=0.55pt, fill=white, circle,
                  inner sep=0pt, minimum size=4.6mm, font=\scriptsize},
  dot/.style   = {ink, fill=ink, circle, inner sep=0pt, minimum size=1.5mm},
  % 글자
  sname/.style = {font=\small, inner sep=1.5pt},
  note/.style  = {font=\scriptsize, soft, inner sep=1.5pt},
  wnote/.style = {font=\scriptsize, warn, inner sep=1.5pt},
  anote/.style = {font=\scriptsize, accent, inner sep=1.5pt},
  axis/.style  = {ink, line width=0.5pt},
  tick/.style  = {font=\scriptsize, soft},
  % 묶음 상자 — 서브시스템 표시
  grp/.style   = {draw=soft, line width=0.4pt, densely dotted, rounded corners=1mm,
                  inner sep=3mm},
}
   (assets/src/ 안에서)
% 빌드:     bash _tools/tikz2svg.sh assets/src/w02-node-topic.tex assets/w02-node-topic.svg
%
% 왜 TikZ 인가
%   손으로 SVG 좌표를 쓰면 화살촉·선 굵기·글자 크기가 그림마다 미묘하게 달라진다.
%   그것이 "자료가 어설퍼 보인다"의 정체다. 스타일을 한 번만 정하고 상대좌표로
%   배치하면 좌표를 고쳐도 그림이 무너지지 않는다.
%   수식은 본문과 같은 TeX 글꼴로 나온다.
%
% 한글
%   ko.TeX(DVI 경로)를 쓴다. XeLaTeX 은 TikZ 그래픽을 PDF special 로 내보내는데
%   dvisvgm 이 그것을 받으려면 Ghostscript < 10.01 이 필요하다. 이 PC 에는
%   10.07 이 깔려 있어 그 경로가 막힌다. latex -> DVI -> dvisvgm 이 유일하게
%   동작하는 조합이다.

\usepackage[hangul]{kotex}
\usepackage{tikz}
\usepackage{amsmath}
\usepackage{xcolor}
\usetikzlibrary{arrows.meta, positioning, calc, fit, backgrounds, decorations.pathmorphing}

% ---- 색 --------------------------------------------------------------------
% 강조는 하나만. 나머지는 검정. 흑백 인쇄에서도 읽혀야 한다.
\definecolor{ink}{HTML}{1A1A1A}
\definecolor{soft}{HTML}{666666}
\definecolor{accent}{HTML}{7C3AED}
\definecolor{warn}{HTML}{B91C1C}

% 계층 색 — Simulink 모델의 블록 색과 같은 규칙을 쓴다
\definecolor{guid}{HTML}{1D4ED8}    % 유도
\definecolor{ctrl}{HTML}{EA580C}    % 제어
\definecolor{thr}{HTML}{D97706}     % 추진기
\definecolor{plant}{HTML}{16A34A}   % 운동모델
\definecolor{nav}{HTML}{7C3AED}     % 항법 · 신호처리
\definecolor{pathc}{HTML}{0D9488}   % 계획 경로 · 항적
\definecolor{errc}{HTML}{C2410C}    % 오차
\definecolor{flow}{HTML}{0369A1}    % 조류 · 바람

% ---- 선과 화살촉 -----------------------------------------------------------
\tikzset{
  sig/.style   = {ink, line width=0.5pt, -{Stealth[length=4pt, width=3pt]}},
  wire/.style  = {ink, line width=0.5pt},
  fb/.style    = {ink, line width=0.5pt, densely dashed,
                  -{Stealth[length=4pt, width=3pt]}},
  asig/.style  = {accent, line width=0.5pt, densely dashed,
                  -{Stealth[length=4pt, width=3pt]}},
  % 블록 — 흰 바탕, 직각, 검은 테두리. 색은 테두리에만 준다
  blk/.style   = {draw=ink, line width=0.55pt, fill=white,
                  minimum width=18mm, minimum height=8.5mm, inner sep=2pt, font=\small},
  gblk/.style  = {blk, draw=guid},
  cblk/.style  = {blk, draw=ctrl},
  tblk/.style  = {blk, draw=thr},
  pblk/.style  = {blk, draw=plant},
  nblk/.style  = {blk, draw=nav},
  % 합산점 · 분기점
  sum/.style   = {draw=ink, line width=0.55pt, fill=white, circle,
                  inner sep=0pt, minimum size=4.6mm, font=\scriptsize},
  dot/.style   = {ink, fill=ink, circle, inner sep=0pt, minimum size=1.5mm},
  % 글자
  sname/.style = {font=\small, inner sep=1.5pt},
  note/.style  = {font=\scriptsize, soft, inner sep=1.5pt},
  wnote/.style = {font=\scriptsize, warn, inner sep=1.5pt},
  anote/.style = {font=\scriptsize, accent, inner sep=1.5pt},
  axis/.style  = {ink, line width=0.5pt},
  tick/.style  = {font=\scriptsize, soft},
  % 묶음 상자 — 서브시스템 표시
  grp/.style   = {draw=soft, line width=0.4pt, densely dotted, rounded corners=1mm,
                  inner sep=3mm},
}
   (assets/src/ 안에서)
% 빌드:     bash _tools/tikz2svg.sh assets/src/w02-node-topic.tex assets/w02-node-topic.svg
%
% 왜 TikZ 인가
%   손으로 SVG 좌표를 쓰면 화살촉·선 굵기·글자 크기가 그림마다 미묘하게 달라진다.
%   그것이 "자료가 어설퍼 보인다"의 정체다. 스타일을 한 번만 정하고 상대좌표로
%   배치하면 좌표를 고쳐도 그림이 무너지지 않는다.
%   수식은 본문과 같은 TeX 글꼴로 나온다.
%
% 한글
%   ko.TeX(DVI 경로)를 쓴다. XeLaTeX 은 TikZ 그래픽을 PDF special 로 내보내는데
%   dvisvgm 이 그것을 받으려면 Ghostscript < 10.01 이 필요하다. 이 PC 에는
%   10.07 이 깔려 있어 그 경로가 막힌다. latex -> DVI -> dvisvgm 이 유일하게
%   동작하는 조합이다.

\usepackage[hangul]{kotex}
\usepackage{tikz}
\usepackage{amsmath}
\usepackage{xcolor}
\usetikzlibrary{arrows.meta, positioning, calc, fit, backgrounds, decorations.pathmorphing}

% ---- 색 --------------------------------------------------------------------
% 강조는 하나만. 나머지는 검정. 흑백 인쇄에서도 읽혀야 한다.
\definecolor{ink}{HTML}{1A1A1A}
\definecolor{soft}{HTML}{666666}
\definecolor{accent}{HTML}{7C3AED}
\definecolor{warn}{HTML}{B91C1C}

% 계층 색 — Simulink 모델의 블록 색과 같은 규칙을 쓴다
\definecolor{guid}{HTML}{1D4ED8}    % 유도
\definecolor{ctrl}{HTML}{EA580C}    % 제어
\definecolor{thr}{HTML}{D97706}     % 추진기
\definecolor{plant}{HTML}{16A34A}   % 운동모델
\definecolor{nav}{HTML}{7C3AED}     % 항법 · 신호처리
\definecolor{pathc}{HTML}{0D9488}   % 계획 경로 · 항적
\definecolor{errc}{HTML}{C2410C}    % 오차
\definecolor{flow}{HTML}{0369A1}    % 조류 · 바람

% ---- 선과 화살촉 -----------------------------------------------------------
\tikzset{
  sig/.style   = {ink, line width=0.5pt, -{Stealth[length=4pt, width=3pt]}},
  wire/.style  = {ink, line width=0.5pt},
  fb/.style    = {ink, line width=0.5pt, densely dashed,
                  -{Stealth[length=4pt, width=3pt]}},
  asig/.style  = {accent, line width=0.5pt, densely dashed,
                  -{Stealth[length=4pt, width=3pt]}},
  % 블록 — 흰 바탕, 직각, 검은 테두리. 색은 테두리에만 준다
  blk/.style   = {draw=ink, line width=0.55pt, fill=white,
                  minimum width=18mm, minimum height=8.5mm, inner sep=2pt, font=\small},
  gblk/.style  = {blk, draw=guid},
  cblk/.style  = {blk, draw=ctrl},
  tblk/.style  = {blk, draw=thr},
  pblk/.style  = {blk, draw=plant},
  nblk/.style  = {blk, draw=nav},
  % 합산점 · 분기점
  sum/.style   = {draw=ink, line width=0.55pt, fill=white, circle,
                  inner sep=0pt, minimum size=4.6mm, font=\scriptsize},
  dot/.style   = {ink, fill=ink, circle, inner sep=0pt, minimum size=1.5mm},
  % 글자
  sname/.style = {font=\small, inner sep=1.5pt},
  note/.style  = {font=\scriptsize, soft, inner sep=1.5pt},
  wnote/.style = {font=\scriptsize, warn, inner sep=1.5pt},
  anote/.style = {font=\scriptsize, accent, inner sep=1.5pt},
  axis/.style  = {ink, line width=0.5pt},
  tick/.style  = {font=\scriptsize, soft},
  % 묶음 상자 — 서브시스템 표시
  grp/.style   = {draw=soft, line width=0.4pt, densely dotted, rounded corners=1mm,
                  inner sep=3mm},
}

\usetikzlibrary{shapes.geometric}

\tikzset{
  %  판정 — 마름모. 글자가 두 줄이라 폭을 넉넉히 준다
  dec/.style = {draw=ink, line width=0.55pt, fill=white, diamond, aspect=2.6,
                inner sep=1pt, font=\scriptsize, align=center,
                minimum width=50mm, minimum height=17mm},
  %  단계 — 네모. 목표가 무엇으로 바뀌는지를 적는다
  pbox/.style = {draw=guid, line width=0.6pt, fill=white, rounded corners=0.8mm,
                 inner sep=3pt, font=\scriptsize, align=center,
                 minimum width=62mm, minimum height=13mm},
  %  시작과 끝 — 둥근 띠
  capsule/.style = {draw=ink, line width=0.6pt, fill=white, rounded corners=3mm,
                inner sep=3pt, font=\small, minimum width=40mm, minimum height=9mm},
  ylab/.style = {font=\scriptsize, ink, inner sep=1.5pt},
  nlab/.style  = {font=\scriptsize, warn, inner sep=1.5pt},
}

\begin{document}
\begin{tikzpicture}[x=1mm, y=1mm]

% ===========================================================================
% 가운데 줄 — 시작에서 끝까지
% ===========================================================================
\node[capsule]  (st) at (0,0) {접안 시작};

\node[pbox] (p1) at (0,-30)
     {\textbf{1단계 · 접근}\\
      목표 위치 = \textbf{접근점} $(x_a,\,y_a)$\\
      목표 선수각 = \textbf{가는 방향}};

\node[dec]  (d1) at (0,-58)
     {접근점까지 남은 거리가\\ $R_{dock} = 4$ m 보다 가까운가?};

\node[pbox] (p2) at (0,-86)
     {\textbf{2단계 · 선수각 정렬}\\
      목표 위치 = 접근점 \textbf{그대로}\\
      목표 선수각 = \textbf{접안 선수각} $\psi_f$};

\node[dec]  (d2) at (0,-114)
     {선수각 오차가\\ $5^\circ$ 보다 작은가?};

\node[pbox] (p3) at (0,-142)
     {\textbf{3단계 · 접안}\\
      목표 위치 = \textbf{접안점} $(x_d,\,y_d)$\\
      목표 선수각 = $\psi_f$ \textbf{유지}};

\node[dec]  (d3) at (0,-170)
     {접안점까지 남은 거리가\\ $R_{stop} = 1.5$ m 보다 가까운가?};

\node[capsule]  (en) at (0,-196) {접안 끝 — \textbf{gate} $= 0$};

% --- 아래로 흐르는 선 -------------------------------------------------------
\draw[sig] (st) -- (p1);
\draw[sig] (p1) -- (d1);
\draw[sig] (d1) -- (p2);   \node[ylab, anchor=west] at (1.6,-72) {예};
\draw[sig] (p2) -- (d2);
\draw[sig] (d2) -- (p3);   \node[ylab, anchor=west] at (1.6,-128) {예};
\draw[sig] (p3) -- (d3);
\draw[sig] (d3) -- (en);   \node[ylab, anchor=west] at (1.6,-184) {예};

% --- 아니오 — 같은 단계를 한 스텝 더 돈다 -----------------------------------
%     되돌아가는 선은 오른쪽 한 줄(x = 56)에 모아 서로 겹치지 않게 한다
\draw[wire, warn] (d1.east) -- (56,-58) -- (56,-30);
\draw[sig,  warn] (56,-30) -- (p1.east);
\node[nlab, anchor=south] at (56,-56.5) {아니오};

\draw[wire, warn] (d2.east) -- (56,-114) -- (56,-86);
\draw[sig,  warn] (56,-86) -- (p2.east);
\node[nlab, anchor=south] at (56,-112.5) {아니오};

\draw[wire, warn] (d3.east) -- (56,-170) -- (56,-142);
\draw[sig,  warn] (56,-142) -- (p3.east);
\node[nlab, anchor=south] at (56,-168.5) {아니오};

\node[nlab, anchor=north west, text width=48mm] at (62,-26)
     {\textbf{아니오}는 되돌아가는 것이 아니다.
      \textbf{같은 단계를 한 스텝 더} 도는 것뿐이다 —
      제어기는 매 스텝 이 그림을 처음부터 다시 읽는다};

% ===========================================================================
% 왼쪽 — 각 단계가 무엇을 하는지, 왜 그렇게 하는지
% ===========================================================================
\node[anchor=north east, font=\scriptsize, text width=74mm, align=left] at (-40,-22)
     {\textbf{왜 선석으로 바로 가지 않는가} \\[2pt]
      선석 앞은 좁다. 비스듬히 들어가면 뱃머리나 고물이 먼저 닿는다.
      \textbf{접근점}은 선석에서 부두 반대쪽으로 $L_{app} = 15$ m 물러난 점이다.
      여기까지는 자세를 따지지 않고 빨리 간다 —
      가는 방향으로 선수를 두면 물 저항이 가장 작다};

\node[anchor=north east, font=\scriptsize, text width=74mm, align=left] at (-40,-78)
     {\textbf{왜 제자리에서 먼저 돌리는가} \\[2pt]
      움직이면서 돌면 배가 호를 그리며 선석 쪽으로 쏠린다.
      접근점에 \textbf{선 채로} 선수각만 맞춰 두면, 남은 일은 똑바로 미는 것
      하나뿐이다. $5^\circ$ 는 이 모델이 쓰는 문턱이다 —
      더 조이면 바람에 흔들려 영영 다음 단계로 못 넘어간다};

\node[anchor=north east, font=\scriptsize, text width=74mm, align=left] at (-40,-134)
     {\textbf{마지막은 선수각을 묶어 둔 채 민다} \\[2pt]
      목표 선수각을 $\psi_f$ 로 \textbf{고정}한다. 위치가 바뀌어도 선수각
      지령은 따라 움직이지 않는다 — 그래야 배가 선석과 나란한 자세를
      유지한 채 붙는다. 이 구간의 횡방향 힘을 \textbf{선수추진기}가 낸다};

\node[anchor=north east, font=\scriptsize, text width=74mm, align=left] at (-40,-190)
     {\textbf{gate 가 내려가면 바깥 고리가 속도 지령을 0 으로 낸다} \\[2pt]
      추력을 끊는 것이 아니라 \textbf{목표 속도를 0 으로} 주는 것이다.
      그래서 남은 속도는 안쪽 속도 고리가 되돌려 세운다 \\[2pt]
      7단(전진)만 이 자리가 다르다 — gate 를 $0 \sim 1$ 의 \textbf{감속 계수}
      $\min(1,\, d/R_{brake})$ 로 두고, 도착 판정에
      $\sqrt{u^2+v^2} < U_{berth}$ 를 \textbf{하나 더} 본다};

% ===========================================================================
% 단계 래치 — 이 그림에서 가장 자주 놓치는 것
% ===========================================================================
\draw[grp, draw=warn] (-40,-178) rectangle (52,-20);
\node[font=\scriptsize\bfseries, warn, anchor=south west] at (-39.5,-19.4)
     {단계는 한 번 올라가면 \textbf{내려오지 않는다} (래치)};

% ===========================================================================
% 아래 띠 — 세 전략이 이 흐름의 어디에서 갈리는가
% ===========================================================================
\draw[wire, soft] (-118,-212) -- (116,-212);

\node[anchor=north west, font=\scriptsize, text width=112mm] at (-118,-217)
     {\textbf{세 전략은 이 흐름을 그대로 쓴다 — 갈리는 자리는 셋뿐이다} \\[3pt]
      $\bullet$ \textbf{접안 선수각 $\psi_f$} — 옆걸음은 $\psi_f = \psi_t$
        (선석 방위와 같게), 전진은 $\psi_f = \psi_t - s\,\dfrac{\pi}{2}$
        (선석 방위에 수직하게). 접근점은 \textbf{두 경우가 같은 점}이고,
        거기서 선수각만 $90^\circ$ 틀면 옆걸음이 전진으로 바뀐다 \\[3pt]
      $\bullet$ \textbf{위치 오차를 어느 축에서 재는가} — 선체고정 갈래는
        오차를 선체 축(앞뒤·좌우)으로 돌려서 재고, NED 갈래는 북·동 축에서
        그대로 재어 속도 지령을 \textbf{벡터로} 만든다. 단계 판정은 둘 다
        똑같이 거리와 각도로 한다 \\[3pt]
      $\bullet$ \textbf{마지막 판정에서 gate 를 끊는가 줄이는가} — 옆걸음과
        NED 는 $R_{stop}$ 안에 들면 gate 를 0 으로 \textbf{끊는다}. 전진은
        빠른 축(전후)으로 들어와 관성으로 더 밀리므로, $R_{brake}$ 안에서
        gate 를 거리에 비례해 \textbf{줄이고} 속도까지 느려진 뒤에 끊는다};

\node[anchor=north west, font=\scriptsize, text width=112mm] at (4,-217)
     {\textbf{단계가 실제로 언제 넘어갔는가} — 모델을 돌려 잰 값 \\[3pt]
      \begin{tabular}{@{}lrrrr@{}}
        & 1$\to$2 & 2$\to$3 & 완료 & 끝 오차 \\[1pt]
        5단 옆걸음 & 48.70 s & 54.65 s & 78.70 s & 0.545 m \\
        7단 전진   & 48.95 s & 57.80 s & 72.90 s & 0.118 m \\
        8단 NED    & 115.45 s & 121.40 s & 145.95 s & 0.554 m \\
      \end{tabular} \\[3pt]
      1$\to$2 가 5단과 7단에서 거의 같은 것은 \textbf{접근점이 같은 점}이기
      때문이다. 갈리는 것은 그 다음부터다 \\[2pt]
      7단의 끝 오차가 유독 작은 것은 거기만 gate 를 끊지 않고 $R_{brake}$
      안에서 \textbf{줄여} 들어오기 때문이다 — 줄이지 않고 끊어만 두면
      gate 가 내려간 뒤에도 남은 속도로 더 나아가 1.713 m 를 남긴다};

\end{tikzpicture}
\end{document}
